\documentclass[11pt,a4paper]{article}
\usepackage[utf8]{inputenc}
\usepackage[T1]{fontenc}
\usepackage[portuguese]{babel}
\usepackage{xcolor}
\usepackage{hyperref}
\usepackage{geometry}
\usepackage{listings}
\usepackage{tcolorbox}
\geometry{margin=2.5cm}
\definecolor{primary}{RGB}{41,128,185}
\definecolor{secondary}{RGB}{44,62,80}
\definecolor{codebg}{RGB}{240,240,240}
\lstset{backgroundcolor=\color{codebg},basicstyle=\ttfamily\small,breaklines=true,frame=single}
\title{\color{secondary}\Huge\textbf{Solução: Exercício 2.2 - Namespaces}}
\author{\color{primary}DevOps com Kubernetes -- 2026}
\date{}
\begin{document}
\maketitle


\textbf{Guia de Solução:}

\subsection*{\color{secondary}Passo a Passo}

1. Crie os namespaces com \texttt{kubectl create namespace dev} e \texttt{kubectl create namespace prod}.
2. Aplique os manifestos usando \texttt{-n dev} e \texttt{-n prod} ou declare \texttt{metadata.namespace}.
3. Compare \texttt{kubectl get svc -n dev} com \texttt{kubectl get svc -n prod}.
4. Teste DNS entre namespaces usando \texttt{nome-do-service.nome-do-namespace}.

\end{document}
